%% ============================================================
%% AULA 4 — Aprendendo a rede: Parte 2 (retropropagação)
%% GBC073 — Inteligência Computacional (FACOM/UFU)
%%
%% Versão sucinta: roteiro do apresentador (poucas palavras,
%% mesmas definições e figuras do aula04.tex)
%%
%% Esqueleto com esboço de conteúdo — a desenvolver.
%% Ementa (ficha SEI 5116656): item 2 — regra delta e backpropagation.
%%
%% Compilar:  xelatex aula04-sucinta.tex (duas passadas)
%% ============================================================
\documentclass[aspectratio=169, 11pt]{beamer}

\makeatletter
\def\input@path{{./}{../}}
\makeatother

\usetheme{InteligenciaComputacional}   % ANTES do babel: fontspec vem primeiro
\usepackage{babel}
\babelprovide[import, main]{brazilian}

\title[Aula 4 — Retropropagação]{Aula 4: Aprendendo a rede — Parte 2: retropropagação}
\subtitle[GBC073 — Inteligência Computacional]{Regra da cadeia no grafo computacional \textbullet\ autograd \textbullet\ PyTorch}
\author[Prof.\ Marcelo Albertini]{Prof.\ Marcelo Keese Albertini}
\institute{GBC073 — Inteligência Computacional \textbullet\ Universidade Federal de Uberlândia}
\date{2026}

\begin{document}

% ------------------------------------------------------------ capa
\begin{frame}[plain, noframenumbering]
  \titlepage
\end{frame}

% ------------------------------------------------------------ roteiro
\begin{frame}{Roteiro da aula}
  \begin{itemize}\setlength{\itemsep}{0.5ex}
    \item<1-> \textbf{O problema do crédito:} quem é responsável por cada erro?
    \item<2-> \textbf{A regra da cadeia no grafo computacional:} o erro viaja para trás
    \item<3-> \textbf{A regra delta generalizada:} a atualização de cada peso
    \item<4-> \textbf{Autograd:} a implementação em PyTorch
  \end{itemize}
\end{frame}

% ------------------------------------------------------------
\section{O problema do crédito}

\begin{frame}{O problema do crédito}
  \begin{itemize}\setlength{\itemsep}{0.4ex}
    \item Rede $=$ composição de camadas: $f = f_L \circ \cdots \circ f_1$
    \item Erro medido na saída $\erro$: culpa \emph{compartilhada} de todos os pesos
    \item Atribuir a cada $\pesoij{j}{i}$ a sua parcela de responsabilidade
    \item Regra do perceptron só alcança a \emph{última} camada — ocultas sem
          sinal de aprendizado
    \item Retropropagação resolve a atribuição de crédito: gradiente camada a camada
  \end{itemize}
\end{frame}

% ------------------------------------------------------------
\section{A regra da cadeia no grafo computacional}

\begin{frame}{A regra da cadeia no grafo computacional}
  \begin{itemize}\setlength{\itemsep}{0.4ex}
    \item Rede como grafo computacional: nós $=$ operações, arestas $=$ tensores
    \item Gradientes locais de cada nó, encadeados pela regra da cadeia
    \item Propaga para trás a \emph{sensibilidade} do erro a cada ativação
    \item Custo da retropropagação: comparável ao da propagação adiante
  \end{itemize}
\end{frame}

% ------------------------------------------------------------
\section{A regra delta generalizada}

\begin{frame}{A regra delta generalizada}
  \begin{itemize}\setlength{\itemsep}{0.4ex}
    \item O erro retropropagado na camada $j$: \(\delta_j = \ativ'(\campo_j) \sum_k \delta_k \pesoij{k}{j}\).
    \item Atualização: \(\Delta \pesoij{j}{i} = \taxa\, \delta_j\, y_i\).
    \item A regra delta (Widrow--Hoff) é o caso de uma camada só.
    \item Funções de ativação diferenciáveis viram requisito — adeus degrau.
  \end{itemize}
\end{frame}

% ------------------------------------------------------------
\section{Autograd e a implementação em PyTorch}

\begin{frame}{Autograd e a implementação em PyTorch}
  \begin{itemize}\setlength{\itemsep}{0.4ex}
    \item Diferenciação automática: grafo computacional registrado durante o
          \emph{forward} (\texttt{tensor.requires\_grad})
    \item \texttt{loss.backward()} percorre o grafo e acumula gradientes
    \item \texttt{optimizer.step()} aplica a atualização; \texttt{zero\_grad()}
          entre passos
    \item Não se escreve retropropagação à mão: escreve-se o \emph{forward} e o
          framework deriva
  \end{itemize}
  \begin{ideiabox}
    Retropropagação $=$ regra da cadeia de trás para frente — o que viaja para
    trás é a sensibilidade do erro a cada peso.
  \end{ideiabox}
\end{frame}

\end{document}
